\begin{abstract}

Fine-tuning can reshape model behavior beyond the source task, but it remains unclear whether structured decision training transfers reusable operations across tasks and how explicit constraints alter this transfer. We investigate these questions using Dou Dizhu, a rule-rich decision environment, and compare ordinary task adaptation with rule-constrained fine-tuning across three open-weight language model checkpoints. Cross-task evaluation shows that ordinary fine-tuning produces strongly model-dependent transfer, whereas constraint-enriched supervision yields a more consistently positive improvement over ordinary adaptation, while remaining selective across reasoning categories. Controlled decision probes further show that rule-constrained adaptation strengthens constraint-induced relative preference shifts, although this probability-level change does not translate into improved greedy-decoding compliance. Causal attention-head ablation on Qwen2.5-7B reveals sparse and heterogeneous changes in functional dependence, with a small subset of heads contributing disproportionately to the measured preference shift. Overall, the results suggest that structured fine-tuning induces selective behavioral transfer and localized mechanistic reorganization rather than a uniform improvement in reasoning ability.

\end{abstract}
